\chapter{构建、测试与复现命令}

本项目已开源，完整代码、测试用例与图表生成脚本见 GitHub 仓库：
\begin{center}
  \url{https://github.com/homer-fang/ftcl}
\end{center}

克隆仓库后，可按下列命令完成构建、测试与实验复现。以下路径以 WSL 下的 \texttt{/mnt/d/ftcl/ftcl} 为例；若从 GitHub 克隆，请将工作目录替换为本地克隆路径。

\section{CPU 构建与测试}
\begin{lstlisting}[language=bash]
git clone https://github.com/homer-fang/ftcl.git
cd ftcl
cmake -S . -B build -DCMAKE_BUILD_TYPE=Release \
  -DBUILD_TESTS=ON \
  -DFTCL_ENABLE_CUDA=OFF \
  -DFTCL_CXX_STANDARD=20
cmake --build build -j8
ctest --test-dir build --output-on-failure
\end{lstlisting}

\section{CUDA 构建}
\begin{lstlisting}[language=bash]
cmake -S . -B /tmp/ftcl-build-geometry-cuda \
  -DCMAKE_BUILD_TYPE=Release \
  -DBUILD_TESTS=ON \
  -DFTCL_ENABLE_CUDA=ON \
  -DFTCL_CXX_STANDARD=20
cmake --build /tmp/ftcl-build-geometry-cuda -j8
\end{lstlisting}

\section{核心基准测试}
\begin{lstlisting}[language=bash]
./build/test/bench_ftcl ./docs/data/benchmarks
python3 ./docs/scripts/benchmarks/plot_benchmarks.py \
  ./docs/data/benchmarks ./docs/figures/benchmarks
\end{lstlisting}

\section{几何图表生成}
\begin{lstlisting}[language=bash]
python3 docs/scripts/geometry/generate_geometry_figures.py \
  --build-dir /tmp/ftcl-build-geometry-cuda

python3 docs/scripts/geometry/generate_geometry_study_addons.py \
  --build-dir /tmp/ftcl-build-geometry-cuda
\end{lstlisting}

\section{论文编译}
\begin{lstlisting}[language=bash]
latexmk -xelatex -interaction=nonstopmode -output-directory=build main.tex
\end{lstlisting}

按当前 \texttt{latexmk} 配置，生成的 PDF 输出为 \texttt{build/main.pdf}。
